\section{Marco teórico y trabajos relacionados}
\subsection{A. Patrón de diseño}
Los patrones de diseño constituyen soluciones ya probadas y estructuradas para problemas recurrentes en el desarrollo de software. Su objetivo principal es evitar la duplicación de código y favorecer su reutilización, aportando al mismo tiempo claridad y modularidad en los sistemas. De acuerdo con \cite{gavilanez2022analisis} afirma que los patrones de diseño permiten mantener una lógica de organización en el código, lo cual facilita la creación de software de calidad, con mayor facilidad de mantenimiento y una mejor comprensión por parte de otros desarrolladores. Además, brindan beneficios como la estandarización de soluciones, la reducción de errores comunes y la optimización de tiempos de desarrollo.

\subsection{B. Metodología}
Una metodología, según \cite{merchan2022comparacion}, puede definirse como el conjunto de procesos, habilidades, procedimientos y documentación de apoyo que guían a los desarrolladores en la implementación de nuevos sistemas de información. Estas metodologías se estructuran en sub-fases que funcionan como una hoja de ruta para la planificación, gestión, control y evaluación de los proyectos. Su propósito es proporcionar un marco de trabajo que oriente la selección de las tecnologías más adecuadas en cada etapa del desarrollo, garantizando así un proceso más organizado, eficiente y alineado con los objetivos del sistema.

\subsection{C. Arquitectura de software}
La arquitectura de software se concibe como la estructura fundamental de un sistema, conformada por sus componentes, las interacciones entre ellos y las propiedades que guían su funcionamiento. Más allá de definir cómo se organiza el código, su importancia radica en que establece un marco de referencia que facilita la escalabilidad, el mantenimiento y la comprensión global del sistema. En este sentido, \cite{vera2022arquitectura} señala que, sin una arquitectura bien definida, el desarrollo pierde efectividad, ya que la falta de organización dificulta la claridad y la evolución del software en el tiempo.

\subsection{D. Programación Orientado a Objetos}
La Programación Orientada a Objetos (POO) constituye un paradigma que busca modelar el mundo real de una manera más cercana a la lógica humana. Según \cite{vera2022arquitectura}, este enfoque introduce principios como el encapsulamiento y la abstracción, los cuales permiten la creación de unidades denominadas clases, favoreciendo la reutilización del código y un mejor entendimiento del mismo. Asimismo, \cite{vera2022arquitectura} sostiene que la POO aporta ventajas significativas al permitir la resolución de problemas a partir de códigos ya funcionales, evitando la duplicidad y promoviendo la eficiencia en el desarrollo de sistemas. Este paradigma involucra la definición de clases, objetos, propiedades y métodos, lo que facilita la implementación y organización de los programas, además de mejorar la mantenibilidad y escalabilidad de los proyectos de software.

\subsection{E. Modelo-Vista-Controlador (MVC)}
El patrón Modelo-Vista-Controlador (MVC) surge con el propósito de reducir el esfuerzo en la programación y estandarizar el diseño de aplicaciones que requieren manejar múltiples vistas sincronizadas con los mismos datos \cite{zambrano_patron}. Esta arquitectura organiza una aplicación en tres componentes principales: el Modelo, encargado de la gestión de datos y la lógica de negocio; la Vista, que define la interfaz de usuario; y el Controlador, que actúa como intermediario entre ambos, gestionando las interacciones y solicitudes del usuario. Esta separación favorece la actualización y el mantenimiento del software de manera más sencilla y en menor tiempo, al permitir trabajar de forma independiente en cada componente \cite{zambrano_patron}. De acuerdo con \cite{gavilanez2022analisis}, entre las principales ventajas del MVC destacan la posibilidad de modificar fácilmente un componente sin afectar los demás, la reutilización del modelo para construir diversas vistas sin necesidad de alterarlo, y la facilidad para implementar pruebas unitarias. En conjunto, estas características hacen del MVC promueve la escalabilidad y la eficiencia en el desarrollo de sistemas de software.

\subsection{F. Scrum}
Scrum es una metodología ágil de gestión de proyectos que busca optimizar los procesos de desarrollo mediante ciclos iterativos e incrementales. Según [2], esta metodología se caracteriza por organizar el trabajo en sprints o iteraciones cortas, en las que se definen objetivos específicos y medibles, lo que permite evaluar avances constantes y realizar ajustes oportunos. Dentro de este marco, los roles como el Scrum Master y el Product Owner son esenciales para garantizar la correcta distribución de tareas, la resolución de problemas y la alineación del equipo hacia los objetivos del proyecto.